\documentclass[11pt,a4paper]{article}
\pdfmapfile{+cm.map}
\pdfmapfile{+cmextra.map}
\pdfmapfile{+symbols.map}
\pdfmapfile{+euler.map}
\pdfmapfile{+latxfont.map}
\pdfmapfile{+rsfs.map}
\usepackage[margin=25mm,headheight=15pt]{geometry}
\usepackage[T1]{fontenc}
\usepackage{lmodern}

\usepackage{amsmath,amssymb,amsthm,amscd,mathtools,mathrsfs}
\usepackage{microtype,booktabs,longtable,array,xcolor,graphicx,enumitem,needspace,fancyhdr}
\usepackage[round,authoryear]{natbib}
\usepackage[unicode,colorlinks=true,linkcolor=blue!40!black,citecolor=blue!40!black,urlcolor=blue!40!black]{hyperref}
\usepackage{xurl}
\usepackage{listings}
\usepackage{tabularx}
\lstset{basicstyle=\ttfamily\small,breaklines=true,columns=fullflexible,keepspaces=true}
\usepackage{tikz}
\usepackage{tikz-cd}
\usetikzlibrary{arrows.meta,calc,patterns,positioning}
\usepackage{aliascnt}
\usepackage[nameinlink,noabbrev]{cleveref}
\allowdisplaybreaks[2]
\setcounter{MaxMatrixCols}{32}
\setlength{\emergencystretch}{3em}
\setlength{\parskip}{2pt}
\newtheorem{theorem}{Theorem}[section]
\newaliascnt{proposition}{theorem}
\newtheorem{proposition}[proposition]{Proposition}
\aliascntresetthe{proposition}
\crefname{proposition}{proposition}{propositions}
\newaliascnt{lemma}{theorem}
\newtheorem{lemma}[lemma]{Lemma}
\aliascntresetthe{lemma}
\crefname{lemma}{lemma}{lemmas}
\newaliascnt{corollary}{theorem}
\newtheorem{corollary}[corollary]{Corollary}
\aliascntresetthe{corollary}
\crefname{corollary}{corollary}{corollarys}
\newaliascnt{claim}{theorem}
\newtheorem{claim}[claim]{Claim}
\aliascntresetthe{claim}
\crefname{claim}{claim}{claims}
\newaliascnt{statement}{theorem}
\newtheorem{statement}[statement]{Statement}
\aliascntresetthe{statement}
\crefname{statement}{statement}{statements}
\theoremstyle{definition}
\newaliascnt{definition}{theorem}
\newtheorem{definition}[definition]{Definition}
\aliascntresetthe{definition}
\crefname{definition}{definition}{definitions}
\newaliascnt{remark}{theorem}
\newtheorem{remark}[remark]{Remark}
\aliascntresetthe{remark}
\crefname{remark}{remark}{remarks}
\newaliascnt{example}{theorem}
\newtheorem{example}[example]{Example}
\aliascntresetthe{example}
\crefname{example}{example}{examples}
\newaliascnt{problem}{theorem}
\newtheorem{problem}[problem]{Open problem}
\aliascntresetthe{problem}
\crefname{problem}{open problem}{open problems}
\newaliascnt{question}{theorem}
\newtheorem{question}[question]{Question}
\aliascntresetthe{question}
\crefname{question}{question}{questions}
\newaliascnt{auditrecord}{theorem}
\newtheorem{auditrecord}[auditrecord]{Verification record}
\aliascntresetthe{auditrecord}
\crefname{auditrecord}{verification record}{verification records}
\newaliascnt{test}{theorem}
\newtheorem{test}[test]{Verification}
\aliascntresetthe{test}
\crefname{test}{verification}{verifications}
\newaliascnt{conjecture}{theorem}
\newtheorem{conjecture}[conjecture]{Conjecture}
\aliascntresetthe{conjecture}
\crefname{conjecture}{conjecture}{conjectures}
\newcommand{\C}{\mathbb C}\newcommand{\R}{\mathbb R}\newcommand{\Z}{\mathbb Z}\newcommand{\Q}{\mathbb Q}\newcommand{\N}{\mathbb N}\newcommand{\F}{\mathbb F}
\newcommand{\CC}{\mathbb C}\newcommand{\RR}{\mathbb R}\newcommand{\ZZ}{\mathbb Z}\newcommand{\QQ}{\mathbb Q}\newcommand{\NN}{\mathbb N}
\newcommand{\PP}{\mathbb P}\newcommand{\Pp}{\mathbb P}\newcommand{\Pone}{\mathbb P^1}\newcommand{\HH}{\mathbb H}
\newcommand{\RP}{\mathbb{RP}}\newcommand{\CP}{\mathbb{CP}}
\newcommand{\OO}{\mathcal O}\newcommand{\cO}{\mathcal O}\newcommand{\cM}{\mathcal M}\newcommand{\cF}{\mathcal F}
\newcommand{\h}{\mathfrak h}\newcommand{\hz}{\mathfrak h_z}\newcommand{\Acon}{\mathfrak A}\newcommand{\wt}[1]{\widehat{#1}}
\newcommand{\red}{\mathrm{red}}\newcommand{\tor}{\mathrm{tor}}\newcommand{\dP}{\mathrm{dP}}
\newcommand{\gh}{\widehat\gamma}\newcommand{\uh}{\widehat u}\newcommand{\wh}{\widehat w}\newcommand{\dhat}{\widehat\delta}\renewcommand{\dh}{\widehat\delta}
\newcommand{\dd}{\mathrm d}\newcommand{\ii}{\mathrm i}\newcommand{\eps}{\varepsilon}\newcommand{\partt}{\mathrm{part}}\newcommand{\Bcirc}{B^\circ}\newcommand{\ddbar}{\partial\bar\partial}
\DeclareMathOperator{\Pic}{Pic}\DeclareMathOperator{\Alb}{Alb}\DeclareMathOperator{\Aut}{Aut}\DeclareMathOperator{\NS}{NS}
\DeclareMathOperator{\Hom}{Hom}\DeclareMathOperator{\End}{End}\DeclareMathOperator{\Ext}{Ext}\DeclareMathOperator{\Tor}{Tor}
\DeclareMathOperator{\Spec}{Spec}\DeclareMathOperator{\Proj}{Proj}\DeclareMathOperator{\Sym}{Sym}\DeclareMathOperator{\Diff}{Diff}
\DeclareMathOperator{\GL}{GL}\DeclareMathOperator{\SL}{SL}\DeclareMathOperator{\PSL}{PSL}\DeclareMathOperator{\Sp}{Sp}
\DeclareMathOperator{\im}{im}\DeclareMathOperator{\coker}{coker}\DeclareMathOperator{\rank}{rank}\DeclareMathOperator{\diag}{diag}
\DeclareMathOperator{\ord}{ord}\DeclareMathOperator{\vol}{vol}\DeclareMathOperator{\pr}{pr}\DeclareMathOperator{\id}{id}
\DeclareMathOperator{\divisor}{div}\DeclareMathOperator{\Tr}{Tr}\DeclareMathOperator{\Nm}{Nm}\DeclareMathOperator{\Sha}{Sha}
\DeclareMathOperator{\Imn}{Im}\DeclareMathOperator{\Ren}{Re}\DeclareMathOperator{\trdeg}{tr.deg}
\newcommand{\Uses}[1]{\par\smallskip\noindent\textit{Uses:} #1.\par\smallskip}
\newcommand{\Verified}{}\newcommand{\AuthorSelfCheck}{}\newcommand{\NeedsCheck}{}\newcommand{\Open}{\textbf{Open.}\ }

\newcommand{\StatementRecord}[4]{\par\smallskip\begingroup\upshape\raggedright\noindent\textbf{Label:} \nolinkurl{#1}.\par\noindent\textbf{Status:} #2.\par\noindent\textbf{Check record:} #3.\par\noindent\textbf{Uses:} #4.\par\endgroup\smallskip}
\newcommand{\RungZero}[1]{\par\noindent\textbf{Result and mechanism.} #1\par\smallskip}
\newcommand{\RungOne}{\par\Needspace{6\baselineskip}\noindent\textbf{Proof at the first rung.}\par\nobreak\smallskip}
\newcommand{\CurrentChecker}{GPT-6 Astra (OpenAI, ChatGPT Pro mode, effort Pro), 18 November 2026}
\newaliascnt{cited}{theorem}
\newtheorem{cited}[cited]{Cited result}
\aliascntresetthe{cited}
\crefname{cited}{cited result}{cited results}
\newaliascnt{refuted}{theorem}
\newtheorem{refuted}[refuted]{Correction}
\aliascntresetthe{refuted}
\crefname{refuted}{correction}{corrections}
\newaliascnt{computed}{theorem}
\newtheorem{computed}[computed]{Computation}
\aliascntresetthe{computed}
\crefname{computed}{computation}{computations}
\newaliascnt{openq}{theorem}
\newtheorem{openq}[openq]{Open question}
\aliascntresetthe{openq}
\crefname{openq}{open question}{open questions}

\providecommand{\rk}{\operatorname{rk}}
\providecommand{\Ann}{\operatorname{Ann}}
\providecommand{\Log}{\operatorname{Log}}
\providecommand{\vir}{\mathrm{vir}}


\newcommand{\ExternalPaperRef}[3]{\ifmmode\text{\upshape #1#2 (Paper~#3)}\else #1#2 (Paper~#3)\fi}
\newcommand{\CompendiumRef}[2]{\ifmmode\text{\upshape S6\_09, Paper~#1, label~\nolinkurl{#2}}\else S6\_09, Paper~#1, label~\nolinkurl{#2}\fi}

\begin{document}
\section{Paper I frontmatter proof}
% Proposed Paper I frontmatter.  This file reorganizes established body
% statements; it does not add verification events to those statements.
% Shared series overview, methodology, navigation and dependency graphic
% are supplied by the series assembler.

\subsection{The main theorems}
\label{series:I:main-theorems}

We use the Claude--Alp\"oge construction with its specified periodic
cusp fan, clockwise finite root generators and affine translations.
Its base is $B=\Pp^1_t$, with special values $0,1,\infty$, and
$B^\circ=B\setminus\{0,1,\infty\}$.  The ordered homology lattice and
period matrix are
\[
 \Lambda=\Z\langle\widehat\gamma,\widehat u,
                  \widehat w,\widehat\delta\rangle,
 \qquad
 \Pi_c=\begin{pmatrix}6\mu&\tau&1&0\\
                       \beta_{\rm part}+c&\mu&0&1\end{pmatrix}.
\]
Write $d=\widehat\delta$, and use $\gamma,u,w,\delta$ for the dual
integral covectors when they occur in lattice forms.  The parameter
$w=e^{-\pi i c}$ and its square $u=w^2$ are distinguished by their
arguments or by the words ``parameter'' and ``covector.''  The
particular solution is fixed by
$(\beta_{\rm part}+\tau)(\infty)=0$; $\mu$ is the cusp-selected
zero of $\wp_\tau$.  A statement about $t=0$ or $t=1$ uses the
specified smooth finite root model.  Unless further data are
displayed, $X_c$ denotes the standard $(3,4)$ sphere member with
$v_1=\varepsilon=(1,2,-4,0)^{\mathsf t}$,
$v_2=-\varepsilon'=(-1,-3,3,0)^{\mathsf t}$ and the untwisted
cusp collar.  We write $e=(0,1)^{\mathsf t}=\Pi_c d$ for the
invariant vertical vector.  The broader admissible translation and higher
multiplicity data are stated in Theorems B and F.

\hypertarget{main-I-A}{}
\paragraph{Theorem A (The period, its exact domain, and its arithmetic norm).}
\label{series:I:A}
The missing period and the normalized norm are
\[
 \beta_{\rm part}=\frac{6i}{\pi}\Log
 \frac{\vartheta_1(\mu,\tau)}{2\sqrt3\,\eta_D(\tau)},
 \qquad
 N(P)=\left|\frac{\vartheta_1(\mu,\tau)}{\eta_D(\tau)}\right|
           e^{-\pi(\Im\mu)^2/\Im\tau},
 \quad P=[\mu].
\]
The displayed logarithm is the branch fixed by the cusp condition.
Putting $L=\log N$ and
$D_c=\Im(\beta_{\rm part}+c)-6(\Im\mu)^2/\Im\tau$, one has
\[
 D_{
 \rm part}=\frac6\pi\log\frac{N}{2\sqrt3},\qquad
 \max N=\sqrt2\quad(t=1),\qquad
 \max D_{\rm part}=-\frac3\pi\log6.
\]
All ordinary fibres and both normalized finite root fibres are
compact complex tori exactly on
\[
 \Im c<H_*:=\frac3\pi\log6,
 \qquad 0<|w|<216.
\]
There is no coarse critical point of $L$ away from $t=1$;
$t=0$ is a noncritical $C^1$ point with $L_t(0)=7/384$.
For every angle $\theta$, $L(1+re^{i\theta})$ decreases strictly
from $\log\sqrt2$ to $-\infty$ as $r$ increases from zero.
Every level below the maximum is one Jordan curve strictly
star-shaped about $1$, real analytic except for a possible
$C^1$ point with nonzero tangent at $t=0$ on the level
$L=\tfrac14\log3$.  Its bounded complement is the strict
superlevel set; these sets form nested discs.

In any of the explicitly transported pointed Weierstrass models,
for $m\ge2$ and $mP\ne O$,
\[
 N(P)^{m^2}=\frac{N(mP)|\Delta_E|^{(m^2-1)/12}}{|\psi_m(P)|}.
\]
For torsion $P$ of order $n$ and $m\equiv\pm1\pmod n$ this gives
\[
 N(P)^{12}=\frac{|\Delta_E|}{|\psi_m(P)|^{12/(m^2-1)}}.
\]
In particular the five specified values are
\[
\begin{array}{c|c|c}
t&\operatorname{ord}(P)&N(P)^{12}\\\hline
0&3&27\\1&2&64\\32/5&4&10\\
64(47-27\sqrt3)/11&6&17+10\sqrt3\\
64(47+27\sqrt3)/11&6&10\sqrt3-17
\end{array}
\]
and the body gives the complete exact division-polynomial lists
for orders four through eight.  For every number field $K$ and
$t\in K\setminus\{0,1\}$, in the arithmetic model
$y^2=x^3-3t^3(t-1)x/4+t^4(t-1)^2/4$ with
$P=(0,t^2(t-1)/2)$,
\[
 \sum_{\sigma:K\hookrightarrow\C}\log N(P_{\sigma(t)})
 =\sum_{v\nmid\infty}n_v\lambda_v(P_t)
       -[K:\Q]\widehat h_O(P_t).
\]
The normalization is $\lambda_\sigma(P)=-\log N(P)$, with complex
embeddings counted separately.  For rational $t$, the positive
algebraic factor $A_t=\exp(\sum_p\lambda_p(P_t))$ gives
$N(P_t)=A_t e^{-\widehat h_O(P_t)}$.  At the eleven nonspecial
rational CM parameters $t=j/1728$ its values are
\[
\begin{array}{r|c@{\qquad}r|c}
j&A_t&j&A_t\\\hline
8000&\sqrt6&-3375&2\\
54000&2^{2/3}3^{1/4}&-32768&\sqrt3\\
287496&2^{3/4}&-884736&1\\
16581375&2&-12288000&3^{5/12}\\
-884736000&1&-147197952000&1\\
-262537412640768000&1&&
\end{array}
\]
These are exact height identities; no transcendence conclusion
at the stated non-torsion points is included.

\textit{Mechanism.} Theta transformation laws fix the additive
constant, division identities control the norm, and the local
height product formula fixes its arithmetic normalization.
\textit{Relation to literature.} The elliptic surface and section
come from \cite[Remark~3.12, p.~19]{AC2026} and
\cite[Definition~2.3]{Engel2026}; the exact coordinate transports
are retained in \ref{I:weierstrass-comparison} and
\ref{fam:src:B.1.1}.  The norm and height conclusions have the
specific scopes of the body results listed here.
\hyperlink{outline-I-A}{Proof in outline}; full statements:
\ref{fam:src:C.1.1}, \ref{fam:src:C.2.3},
\ref{ari:C.4.2}, \ref{ari:C.4.3}, \ref{ari:C.4.4},
\ref{ari:C.6.1}, \ref{ari:C.6.3}, \ref{ari:C.8.15},
\ref{push:src:D.4.8}.

\hypertarget{main-I-B}{}
\paragraph{Theorem B (All discrete data, the parameter, and real forms).}
\label{series:I:B}
For the fixed period solution and cusp fan, take
\[
 w_1=r_1\varepsilon+k_1d,\qquad
 w_2=r_2\varepsilon'+k_2d,\qquad
 3\nmid r_1,\quad2\nmid r_2,
 \quad p=12l_0-4r_1-3r_2,
\]
where $l_0$ is the first integral winding of the full holomorphic
cusp collar in the stated positive convention.  Two such members
are biholomorphic exactly when
\[
 n=c'-c\in\Z,\qquad p'=ap\ (a\in\{1,-1\}),\qquad
 n\equiv k_2-k'_2\pmod2.
\]
Every biholomorphism is over $B$, and the map with sign $a$ has
fibre-linear part $aI_2$.  The complete invariants are
\[
 \bigl(|p|,e^{-\pi i(c+k_2)}\bigr),\qquad
 |p|\ge1,\quad\gcd(|p|,6)=1,
\]
and every displayed integer occurs.  The additional analytic
collar part and its second integral winding contribute no further
invariant.  For sphere members, $|p|=1$, a fixed datum has period
subgroup $2\Z$.  At each fixed admissible $u=w^2$ there are
exactly two sphere isomorphism classes, exchanged by a positive
circuit of $u$ about zero.

For positive $k,l$ with $\gcd(3k,4l)=1$, put $m_1=3k,m_2=4l$.
Use Viaclovsky's sphere data
$w_j=r_j\varepsilon_j+s_jd$, $\gcd(r_j,m_j)=1$,
$\varepsilon_1=\varepsilon$, $\varepsilon_2=\varepsilon'$, and
$p=m_1m_2l_0-m_2r_1-m_1r_2=\pm1$.
Two members are biholomorphic precisely when their multiplicity
pairs agree and, for some $a=\pm1$ and integer $n=c'-c$,
\[
 p'=ap,\quad r'_j\equiv ar_j\pmod{m_j},\quad
 nr'_1\equiv as_1-s'_1\pmod{k},\quad
 nr'_2\equiv as_2-s'_2\pmod{2l}.
\]
The fixed-data period subgroup is $2kl\Z$.  At fixed admissible
$u$ there are exactly $2kl$ classes for this multiplicity pair,
and a parameter circuit permutes them in one cycle.

For the standard family the reduced finite fibres have Albanese
nomes
\[
 q_2=iw/216,\qquad q_1=-w^2/48^3=(27/64)q_2^2.
\]
An antiholomorphic diffeomorphism $X_c\to X_{c'}$ exists exactly
when $c'-(1-\bar c)\in2\Z$, equivalently $w'=-\bar w$.
The members with real structures are precisely
$\Re c\in\tfrac12+\Z$.  On each,
\[
 \operatorname{Aut}^{\pm}(X_c)=\C^*\rtimes\Z/2,
 \qquad z\longmapsto1/\bar z,
\]
and there are two holomorphic-conjugacy classes of antiholomorphic
involutions.  For negative imaginary $w$ both fixed manifolds
are $S^3$, with the semifree circle action having the circle
$D_2^{\sigma_c}$ as fixed set.  For positive imaginary $w$ one
class is fixed-point free and the other has fixed manifold $S^3$
with the free Hopf circle action, up to the stated circle and
orientation conventions.

\textit{Mechanism.} Integral centralizers and finite norm maps
retain the cusp winding and the parity obstruction; the marked
reflection and cap-meridian calculation determine the real forms.
\textit{Relation to literature.} The construction is that of
\cite{AC2026,Engel2026}; the higher input is Viaclovsky's specified
root construction.  The standard exact $2\Z$ classification is
also Theorem~1.4 of \cite{DuEtAl2026v1}; the version and parameter
comparison is retained in \ref{fam:src:A.5.1} and the credited
appendix, including the constant translation between particular
period solutions.
\hyperlink{outline-I-B}{Proof in outline}; full statements:
\ref{fam:src:B.2.1}, \ref{fam:src:B.2.2}, \ref{fam:src:B.2.3},
\ref{fam:src:B.3.1}, \ref{ari:C.2.1}, \ref{fam:src:B.4.2},
\ref{mapsI:src:A.3.1}, \ref{mapsI:src:A.3.2},
\ref{mapsI:src:A.3.3}.

\hypertarget{main-I-C}{}
\paragraph{Theorem C (The monodromy quotient and every fibre's algebraic dimension).}
\label{series:I:C}
Let $\psi=6\gamma\wedge\delta+u\wedge w$ on $\Lambda$, and
let $P_{d,\Z}$ fix $d$ as a vector in its integral symplectic
group.  Write $J_2=\left(\begin{smallmatrix}0&1\\-1&0\end{smallmatrix}\right)$.
Every member has unique block coordinates
\[
 g(B,x,t)=\begin{pmatrix}1&0&0\\6x&B&0\\
                   t&-x^{\mathsf t}J_2B&1\end{pmatrix},
 \quad B\in\operatorname{SL}_2(\Z),\ x\in\Z^2,\ t\in\Z.
\]
For $S=\left(\begin{smallmatrix}0&-1\\1&0\end{smallmatrix}\right)$,
$T=\left(\begin{smallmatrix}1&1\\0&1\end{smallmatrix}\right)$,
let $\alpha:\operatorname{SL}_2(\Z)\to\Z/12$ be the homomorphism
with $\alpha(S)=9$, $\alpha(T)=1$.  The actual monodromy group is
\[
 \Gamma=\{g(B,x,t):t+6(x_1+x_2+x_1x_2)-\alpha(B)=0\pmod{12}\}.
\]
It has index 12 in $P_{d,\Z}$ and index 24 in $P_\Z$, the group preserving
$\Z d\subset\ker\gamma\subset\Lambda$.  In the real parabolic,
a fundamental set is given by the classical modular region for
$\tau$, the Iwasawa angle $0\le\theta<\pi$,
$0\le x_1,x_2<1$ and $0\le t<12$, with the actual modular
lifts $g(S,0,9)$, $g(T,0,1)$, $g(-I,0,6)$ and the Heisenberg
face identifications
\[
 (x_1,x_2,t)\mapsto(x_1+1,x_2,t-6-6x_2),\quad
 (x_1,x_2,t)\mapsto(x_1,x_2+1,t-6+6x_1).
\]

On $\mathcal D=\operatorname{Sp}_4(\R)/\operatorname U(1,1)$ put
$q_d(J)=\psi(d,Jd)$.  The largest $\Gamma$-invariant open domain
of proper discontinuity is $\mathcal O=\{q_d\ne0\}$;
the same assertion holds for $P_\Z$-invariant open domains and
its action.  The null locus has a dense set of
points with infinite monodromy stabilizer.  Each sign region has
a Hausdorff complex analytic orbifold quotient.  In the family
chart, $\mathcal O_+$ has $\Im\tau>0$, $D_c<0$ and
$q_d=-6/D_c$.  For every admissible $c$ the period map is
injective and immersive on $B^\circ$ and extends to a proper
holomorphic embedding
\[
 \C=B\setminus\{\infty\}\lhook\joinrel\longrightarrow
 \Gamma\backslash\mathcal O_+.
\]

For every smooth fibre its N\'eron--Severi group is the integral
kernel in $\Z^6$ of
\[
 E_{12}-E_{13}\tau+(6E_{23}-E_{14})\mu+E_{24}\beta_c
             +E_{34}(6\mu^2-\tau\beta_c)=0.
\]
With its complex orientation,
$E^2=-2\operatorname{Pf}(E)$, where
$\operatorname{Pf}(E)=E_{12}E_{34}-E_{13}E_{24}+E_{14}E_{23}$.
Let $\mathcal E_c$ be the image on $B^\circ$ of the union of
these zero sets for $E\ne0$ with $\operatorname{Pf}(E)=0$,
and let $\mathcal A_c$ use $\operatorname{Pf}(E)<0$.
They are exactly the elliptic-quotient locus and abelian-surface
locus, respectively, and
\[
 a(F_t)=2\ (t\in\mathcal A_c),\quad
 a(F_t)=1\ (t\in\mathcal E_c\setminus\mathcal A_c),\quad
 a(F_t)=0\ (t\notin\mathcal E_c\cup\mathcal A_c).
\]
These countable unions classify all quotient maps, including
those that mix the two displayed coordinates.  The invariant
class $\eta=6\gamma\wedge\delta+u\wedge w$ has square $-12$;
a fibre with $\operatorname{NS}=\Z\eta$ has no compact curves.

\textit{Mechanism.} The exact congruence kernel controls the real
parabolic quotient, while the integral $(1,1)$ equation and its
Hermitian form determine every fibre stratum with the actual
orientation.
\textit{Relation to literature.} The proper-action question is
posed in \cite[Remark~8.4, pp.~23--24]{Engel2026}.  The stated
answer retains the precise monodromy subgroup and full invariant
open set; it does not rely on replacing that subgroup by its
rational or real closure.
\hyperlink{outline-I-C}{Proof in outline}; full statements:
\ref{fam:src:F.1.1}, \ref{I:monodromy-fundamental-set},
\ref{fam:src:F.2.2}, \ref{fam:src:F.2.3}, \ref{fam:src:F.3.1},
\ref{fam:src:D.1.4}, \ref{fam:src:D.1.6}.

\hypertarget{main-I-D}{}
\paragraph{Theorem D (Cohomology and the ordinary relative differential modules).}
\label{series:I:D}
Every admissible standard sphere member has the following
Dolbeault, Bott--Chern and Aeppli dimensions, with columns
$p=0,1,2,3$ and rows $q=3,2,1,0$:
\[
\begin{array}{c|rrrr}q\backslash p&0&1&2&3\\\hline
3&0&0&0&1\\2&0&1&2&1\\1&1&2&1&0\\0&1&0&0&0
\end{array}_{\bar\partial}\quad
\begin{array}{c|rrrr}q\backslash p&0&1&2&3\\\hline
3&0&0&1&1\\2&0&1&2&1\\1&0&2&1&0\\0&1&0&0&0
\end{array}_{BC}\quad
\begin{array}{c|rrrr}q\backslash p&0&1&2&3\\\hline
3&0&0&0&1\\2&0&1&2&0\\1&1&2&1&0\\0&1&1&0&0
\end{array}_{A}.
\]
The Fr\"olicher sequence satisfies $E_1\ne E_2=E_\infty$;
the only nonzero $d_1$ have rank one at
$(p,q)=(0,1),(1,1),(1,2),(2,2)$.  The Bott--Chern algebra is
$\C\cdot1\oplus\mathcal N$ with $\mathcal N^2=0$.

For the ordinary relative differentials the local presentations
retain both the degree and the full fibre multiplicity.  For
$R=\C\{x,y,z\}$ and $t=xy$,
\[
 \Omega^1_{R/\C\{t\}}=R^3/R(y,x,0)\simeq(x,y)\oplus R,
 \quad\Omega^2\simeq(x,y)\oplus R/(x,y),
 \quad\Omega^3\simeq R/(x,y).
\]
For $t=xyz$, $\Omega^1=R^3/R(yz,xz,xy)$ is torsion-free of
rank two and not locally free on the double locus.  With
$I=(xy,xz,yz)$,
\[
 0\longrightarrow
 \frac{Ru\oplus Rv}{(xu,\,y(u-v),\,zv)}
 \longrightarrow\Omega^2\xrightarrow{\wedge dt}I
 \longrightarrow0,
 \qquad \Omega^3=R/I.
\]
The left term is exactly the torsion and is annihilated by $I$.
For $t=s^m$ in $R=\C\{s,y_1,y_2\}$, the three modules are
\[
 R^2\oplus R/(s^{m-1})\,ds,\qquad
 R\,dy_1\wedge dy_2\oplus(R/(s^{m-1}))^2,
 \qquad R/(s^{m-1})\,ds\wedge dy_1\wedge dy_2.
\]
Consequently the degree-one torsion of $\Omega^1_{X/B}$ is
supported exactly on $S_1\cup S_2$, while degree two also has
torsion along the cusp double locus.

\textit{Mechanism.} Logarithmic residues and root characters
compute the direct images; the evaluated connecting scalar fixes
the remaining Hodge integer.  The ordinary modules follow from
the exact quotient by $dt$ and its exterior powers.
\textit{Relation to literature.} The holomorphic-form vanishing
and Hodge pattern with one undetermined integer are credited to
\cite[Theorem~9.1(vi), p.~65; Proposition~9.20, p.~80]{AC2026}.
The local presentations correct the degree-one torsion assertion
in its Remark~9.21, p.~81, and qualify Remark~1.3, p.~5, by
degree.  The explicit smooth global coefficient atlas remains the
question in \ref{inv:C3.8}.
\hyperlink{outline-I-D}{Proof in outline}; full statements:
\ref{inv:A2.2}, \ref{inv:A3.1}, \ref{inv:C1.1},
\ref{inv:C2.1}, \ref{inv:C3.3}, \ref{inv:A4.1},
\ref{inv:A4.2}.

\hypertarget{main-I-E}{}
\paragraph{Theorem E (Every Picard twist and the deformation parameter).}
\label{series:I:E}
Use the primitive coordinate $M_z=\exp(z\mathfrak g)$ on
$\operatorname{Pic}(X)=\C$, so
$M_n=\mathcal O(n(S_1-S_2))$ for $n\in\Z$.
Let $r_1(n)\in\{0,1,2\}$ and $r_2(n)\in\{0,1,2,3\}$ be the
residues of $n$ modulo 3 and $3n$ modulo 4, respectively, and set
\[
 d(n)=\frac{n-4r_1(n)-3r_2(n)}{12},\quad
 U(n)=\max\{d(n)+1,0\},\quad V(n)=\max\{-d(n)-1,0\},
 \quad\epsilon_n=\begin{cases}0&3\mid n,\\1&3\nmid n.\end{cases}
\]
Then
\[
\begin{aligned}
h^0(T_X\otimes M_n)&=U(n)+U(n-1),\\
h^1(T_X\otimes M_n)&=V(n)+V(n-1)+U(n)+U(n-1)+U(n+2),\\
h^2(T_X\otimes M_n)&=V(n)+V(n-1)+V(n+2)+U(n+6)+\epsilon_n,\\
h^3(T_X\otimes M_n)&=V(n+6).
\end{aligned}
\]
For $z\notin\Z$, let $\mathcal R_z$ be the set of distinct
regular base points for which integers $n,p,q,r$ satisfy
\[
 q\tau+r\mu=p,\qquad z=n-6q\mu-r\beta_c.
\]
It is finite; write $N_z=|\mathcal R_z|$.  Put
$\rho=e^{\pi i/3}$, $\beta_j=\beta_c(z_j)$ at the specified
order-three and order-four marked root centers.  Define
$\ell_1(z)=1+\epsilon_n$ at
$z=n+k(3\beta_1+4-2\rho)$, $k\ne0$, and zero elsewhere;
define $\ell_2(z)=1$ at
$z=n+k(2\beta_2+3-3i)$, $k\ne0$, $n\equiv1,2\pmod4$,
and zero elsewhere.  With
$J_z=2N_z+\ell_1(z)+\ell_2(z)$,
\[
 (h^0,h^1,h^2,h^3)(T_X\otimes M_z)=(0,J_z,1+J_z,0).
\]
The actual cohomology jump loci are locally discrete in this
Picard line, and the generic tuple is $(0,0,1,0)$.

The untwisted tangent tuple is $(1,1,1,0)$, the holomorphic
automorphism group of the standard member is $\C^*$, and its
Kuranishi germ is a smooth reduced line.  For every sphere member
with coprime multiplicities $(3k,4l)$ and the data of Theorem B,
put $z_0=kl$, $d_0=3k+4l$.  On the analytic weighted projective root stack
$\mathfrak B_{k,l}=\Pp(3k,4l)$ the actual root lift satisfies
\[
 \mathcal E=\mathcal O\oplus\mathcal O(-z_0),\quad
 R^q\bar f_*\mathcal O_X=
 \mathcal O,\ \mathcal O\oplus\mathcal O(-z_0),\ \mathcal O(-z_0),
 \quad K_X=\bar f^*\mathcal O(z_0-d_0).
\]
Its tangent tuple is again $(1,1,1,0)$ and its Kuranishi germ is
a smooth reduced line.  For the groupoid derivative
$\kappa(\Phi)=(\partial_c\Phi)\Phi^{-1}$ the classes of
$\partial_c$ and $\partial_u$ are
$e\otimes[\gamma]$ and $-(2\pi i u)^{-1}e\otimes[\gamma]$,
respectively, and are nonzero.  Every stated fixed-data period
family is locally miniversal.

\textit{Mechanism.} Root-line section counts and the exact local
Koszul modules calculate every Picard twist.  The nonzero
Kodaira--Spencer class supplies the actual coordinate on the
one-dimensional Kuranishi germ.
\textit{Relation to literature.} The standard tangent, automorphism
and Kuranishi conclusions are credited to
\cite[Theorems~1.3--1.4, p.~3; Theorem~6.12, pp.~81--82;
Theorem~6.16, p.~85]{DuEtAl2026v1}.  The version-specific
re-derivations and parameter transport are preserved; version~2's
omission of the deformation section changes none of that credit.
\hyperlink{outline-I-E}{Proof in outline}; full statements:
\ref{inv:D4.1}, \ref{inv:D4.2}, \ref{inv:D4.3},
\ref{inv:D4.4}, \ref{inv:D4.5}, \ref{fam:st:A.5.2},
\ref{fam:src:B.4.1}, \ref{push:src:E.2.1}.

\hypertarget{main-I-F}{}
\paragraph{Theorem F (The puncture and the obstruction to every proper filling).}
\label{series:I:F}
For an admissible $(3,4)$ logarithmic datum put
$b=4l_1+3l_2$, $p=12l_0-b$, $a=-p$, with
$3\nmid l_1$ and $l_2$ odd.  After $w=r^6$ there is a proper
holomorphic partial filling over
\[
 (\Pp^1\times\Delta_r)\setminus\{(\infty,0)\}
\]
with smooth total space and reduced normal-crossing parameter
central divisor.  Over a regular finite base point the central
fibre is a cycle of twelve ruled surfaces, with component cover
\[
 y^{12}=t^{4l_1}(t-1)^{3l_2}
\]
and rotations $4l_1,3l_2,-b$ modulo 12.
On $q=Q^{12m}$ and $r=R^n$, the difference between the actual
cusp-extending and parameter-extending origins has deck
differences
\[
 am\gamma-6am^2\delta,\qquad
 an\gamma-6an(n-1)\delta,
\]
so the radial class modulo $\Z\widehat\delta$ is
$-pn\widehat\gamma$ in these fixed conventions.

The punctured family has no proper Hausdorff complex-analytic
filling at $w=0$ after any finite parameter cover.  The conclusion
permits arbitrary singularities, multiplicities and extra central
components, and holds for every admissible datum, including every
sphere member.  It also holds for every admissible datum of the
coprime $(3k,4l)$ construction.  In that case
\[
 N=12kl,\quad b=4l\,l_1+3k\,l_2,\quad a=b-Nl_0,
 \qquad \gcd(l_1,3k)=\gcd(l_2,4l)=1,
\]
the corresponding partial model has $N$ components, and the two
displayed deck differences have $N/2$ in place of 6.

For the untwisted elliptic-surface quotient $X(u)$, whose
logarithmic translations are zero, the actual origins have zero
corner winding.  That family admits a proper flat Hausdorff
analytic filling at $u=0$, obtained by the two-projective-line-bundle
plumbing, proper modification and finite quotient in the body.
The particular one-component central model proposed in
\cite[Conjecture~8.2, p.~22]{Engel2026} cannot be the central
fibre of a proper flat smoothing to these $X(u)$: its virtual
Euler characteristic is zero, whereas $e(X(u))=2$.

\textit{Mechanism.} The intrinsic canonical system recovers the
marked base over a proposed filling.  Finite-orbit vanishing kills
horizontal meridians, while commuting meridians propagate their
ratio across the connected compact corner to the nonzero cusp
coefficient $a$.
\textit{Relation to literature.} The comparison retains Engel's
conjectural status, its specified central model, and the actual
maps between the elliptic-surface and period descriptions.  The
higher family is Viaclovsky's construction; the obstruction is
applied to its precise primitive logarithmic data.
\hyperlink{outline-I-F}{Proof in outline}; full statements:
\ref{night:P2.2}, \ref{night:P2.3}, \ref{night:P2.4},
\ref{night:P2.5}, \ref{night:P2.6}, \ref{night:P2.7},
\ref{night:P2.9}, \ref{night:P2.10}, \ref{night:P2.11}.

\hypertarget{main-I-G}{}
\paragraph{Theorem G (The outer boundary and all marked additive limits).}
\label{series:I:G}
For every $w_0$ with $|w_0|=216$ and every connected open
neighbourhood $U$ of $w_0$, there is no flat proper holomorphic
map from a reduced complex space to $U$ whose fibre at each
$w\in U\cap\{0<|w|<216\}$ is biholomorphic to $X_w$.
The intrinsic divisor $2S_2$ and its Albanese nome $iw/216$
retain the invariant $J(iw/216)$, where
$J(q)=q^{-1}+744+196884q+\cdots$ has the unit circle as a
natural boundary.

Fix $c_*$ with $\Im c_*=H_*$, put $a=\Re c_*+3/2$, and use the
specified order-four root coordinate $h$ at $t=1$ and its fixed
complex-linear fibre coordinates $(\eta,\xi)$.  In the transformed
integral period basis the four vectors are
\[
 (1,1/2),\quad(\tau,1/2+h),\quad(0,1),
 \quad(6h,a+B(h)+c-c_*),
 \qquad B(h)=-3ih^2+O(h^4).
\]
For the total root quotient at fixed $c=c_*$, two central points
are nonseparated exactly when
\[
 \Re(\xi'-\xi)\in G_a=\tfrac12\Z+a\Z\quad(a\in\Q);
\]
every central pair is nonseparated for $a\notin\Q$.  Distinct
noncentral points and distinct base fibres are separated.
The Hausdorff reflection keeps the noncentral complex tori and
has central fibre $\R/G_a\simeq S^1$ for rational $a$, and one
point for irrational $a$.  For reduced $a=m/n$, $n>0$,
$G_a=\operatorname{lcm}(2,n)^{-1}\Z$.

For sequences $c_n\to c_*$ with $\Im c_n<H_*$ and $h_n\to0$,
write
\[
 \delta_n=-D_{c_n}(h_n)>0,\quad
 V_n=\delta_n^{-1}(6\Re h_n,6\Im h_n,\Re(c_n-c_*)),
\]
\[
 L_a=\{(p,q,r):r\in\Z,\ p+r/2\in\Z,
                      \ q+r/2\in\Z,\ ar\in\Z\}.
\]
It has rank three for rational $a$, and equals
$\Z^2\times\{0\}$ for irrational $a$.  Each convergent
subsequence of the closed period lattices has limit
\[
 \Gamma_\infty=\{(\eta,\xi):p\Re\eta+q\Im\eta+r\Re\xi
       +\ell(p,q,r)\Im\xi\in\Z\text{ for every }(p,q,r)\in K\},
\]
where $K\subset L_a$ is primitive,
$K=\{\lambda:\lambda\cdot V_n\text{ has a finite limit}\}$,
and $\ell(\lambda)=\lim\lambda\cdot V_n$.
For labels outside $K$ the absolute values tend to infinity.
Conversely every such primitive $K$ and homomorphism
$\ell:K\to\R$ occurs.  The resulting pointed additive quotient
limits are real tori of every dimension zero through three for
rational $a$, and zero through two for irrational $a$.

\textit{Mechanism.} The Albanese invariant detects the outer
natural boundary.  At a boundary root chart, the exact dual
period matrix and its quadratic cancellation describe every
surviving integral covector, hence both nonseparation and all
closed-group limits.
\textit{Relation to literature.} The elliptic-surface description
is transported through the marked quotient dictionary; the body
states the precise Pontryagin--Chabauty duality input used for
closed-group convergence.  The Hausdorff reflection and additive
limits retain their stated topological categories and maps.
\hyperlink{outline-I-G}{Proof in outline}; full statements:
\ref{push:src:D.1.3}, \ref{push:src:D.2.1},
\ref{push:src:D.2.3}, \ref{push:src:D.3.1}.

\subsection{Proofs in outline}
\label{series:I:outlines}

\hypertarget{outline-I-A}{}
\paragraph{Theorem A.}
\begin{enumerate}[label=\arabic*.]
\item\hypertarget{step-I-A-1}{} The pointed maps
$(x,y)\mapsto(2x,-2\sqrt2y)$ and
$(x_0,y_0)\mapsto(-x_0,y_0/(2i))$ transport the chosen section,
invariant differential and period lattice.  The quotient
exponential identifies the period model with the elliptic-surface
line-bundle quotient
(\hyperref[I:weierstrass-comparison]{\ref*{I:weierstrass-comparison}},
\hyperref[fam:src:B.1.1]{\ref*{fam:src:B.1.1}}).
\item\hypertarget{step-I-A-2}{} The two theta transformation
factors give the prescribed affine period laws; the cusp product
fixes the additive constant and hence the displayed logarithm
(\hyperref[fam:src:C.1.1]{\ref*{fam:src:C.1.1}}).
\item\hypertarget{step-I-A-3}{} The distribution identity gives
the exact division-polynomial norm formula, and $mP=\pm P$
cancels the remaining norm.  The stated polynomial roots and
values then determine every order-four-through-eight
specialization
(\hyperref[ari:C.4.2]{\ref*{ari:C.4.2}},
\hyperref[ari:C.4.3]{\ref*{ari:C.4.3}},
\hyperref[ari:C.4.4]{\ref*{ari:C.4.4}}).
\item\hypertarget{step-I-A-4}{} Exact polynomial bounds away
from the square root point and the controlled local theta
expansion give $\max N=\sqrt2$.  The determinant identity and
its cusp limit then give exactly $\Im c<H_*$
(\hyperref[ari:C.8.1]{\ref*{ari:C.8.1}},
\hyperref[ari:C.8.2]{\ref*{ari:C.8.2}},
\hyperref[ari:C.8.3]{\ref*{ari:C.8.3}},
\hyperref[ari:C.8.4]{\ref*{ari:C.8.4}},
\hyperref[fam:src:C.2.3]{\ref*{fam:src:C.2.3}}).
\item\hypertarget{step-I-A-5}{} The normalized critical factor
satisfies $|1-X|<1$ and
$(t-1)L_t=X(X-2)/24$.  Its radial sign, the coarse derivative
at zero and the cusp limit give the complete critical and level
classification
(\hyperref[ari:C.8.15]{\ref*{ari:C.8.15}},
\hyperref[push:src:D.4.8]{\ref*{push:src:D.4.8}}).
\item\hypertarget{step-I-A-6}{} The local multiplication identity
evaluates every finite-place contribution on the stated minimal
models.  The product formula identifies the archimedean term and
gives the eleven exact CM factors
(\hyperref[ari:C.6.1]{\ref*{ari:C.6.1}},
\hyperref[ari:C.6.3]{\ref*{ari:C.6.3}}).
\end{enumerate}
\hyperlink{main-I-A}{Return to Theorem A}.

\hypertarget{outline-I-B}{}
\paragraph{Theorem B.}
\begin{enumerate}[label=\arabic*.]
\item\hypertarget{step-I-B-1}{} The intrinsic canonical pencil
fixes the base.  Its period centralizer forces $c'-c\in\Z$ and
the fibre-linear part $aI_2$
(\hyperref[fam:src:A.1.1]{\ref*{fam:src:A.1.1}},
\hyperref[fam:src:A.5.1]{\ref*{fam:src:A.5.1}}).
\item\hypertarget{step-I-B-2}{} Solving the two finite norm
equations gives the parity condition; comparing the actual cusp
collar gives $p'=ap$.  The holomorphic correction solves the
remaining gluing, so these necessary equations are sufficient
(\hyperref[fam:src:B.2.1]{\ref*{fam:src:B.2.1}},
\hyperref[fam:src:B.2.2]{\ref*{fam:src:B.2.2}},
\hyperref[fam:src:B.2.3]{\ref*{fam:src:B.2.3}}).
\item\hypertarget{step-I-B-3}{} The higher root pullbacks multiply
the norm images by $k,l$, yielding exactly the congruences modulo
$k,2l$.  Their intersection is $2kl\Z$, and a positive $u$
circuit is the cyclic parameter translation
(\hyperref[fam:src:B.3.1]{\ref*{fam:src:B.3.1}},
\hyperref[push:src:E.1.2]{\ref*{push:src:E.1.2}}).
\item\hypertarget{step-I-B-4}{} The invariant-coordinate lattice,
including its affine half-period, and the two exact finite theta
values give the Albanese nomes
(\hyperref[ari:C.2.1]{\ref*{ari:C.2.1}}).
\item\hypertarget{step-I-B-5}{} The explicit reflected period
and affine lattice maps give $c'\equiv1-\bar c\pmod{2\Z}$.
Their square and holomorphic conjugacy identify the two
involution classes
(\hyperref[fam:src:B.4.2]{\ref*{fam:src:B.4.2}},
\hyperref[mapsI:src:A.3.1]{\ref*{mapsI:src:A.3.1}}).
\item\hypertarget{step-I-B-6}{} The fixed-lift congruences and
cusp units determine the nonempty pieces.  The exact cap-meridian
matrices close them to $S^3$ and determine the semifree or Hopf
circle action
(\hyperref[mapsI:src:A.3.2]{\ref*{mapsI:src:A.3.2}},
\hyperref[mapsI:src:A.3.3]{\ref*{mapsI:src:A.3.3}}).
\end{enumerate}
\hyperlink{main-I-B}{Return to Theorem B}.

\hypertarget{outline-I-C}{}
\paragraph{Theorem C.}
\begin{enumerate}[label=\arabic*.]
\item\hypertarget{step-I-C-1}{} The parabolic multiplication law
and the polarization of $x_1+x_2+x_1x_2$ prove the character law.
The commutator and its conjugates generate its integral kernel,
giving the two exact indices
(\hyperref[fam:src:F.1.1]{\ref*{fam:src:F.1.1}}).
\item\hypertarget{step-I-C-2}{} Actual modular lifts first reduce
the Iwasawa factor, integral Heisenberg translations reduce $x$,
and the center reduces $t$.  Interior uniqueness and the displayed
face maps give the fundamental set
(\hyperref[I:monodromy-fundamental-set]{\ref*{I:monodromy-fundamental-set}}).
\item\hypertarget{step-I-C-3}{} Real parabolic stabilizers are
compact on $q_d\ne0$.  Rational parabolic conjugates of the
explicit null unipotent give a dense infinite-stabilizer locus,
which proves maximality
(\hyperref[fam:src:F.2.2]{\ref*{fam:src:F.2.2}},
\hyperref[fam:src:F.2.3]{\ref*{fam:src:F.2.3}}).
\item\hypertarget{step-I-C-4}{} The modular function $j/1728$
is the period map's left inverse, while the finite root charts
extend it at zero and one.  Its control on compact inverse
images proves properness on the finite $t$-line
(\hyperref[fam:src:F.3.1]{\ref*{fam:src:F.3.1}}).
\item\hypertarget{step-I-C-5}{} Wedge the period rows with an
integral alternating form to obtain the exact $(1,1)$ equation.
Its Hermitian determinant and the complex orientation give
$E^2=-2\operatorname{Pf}(E)$ and the quotient/projectivity tests
(\hyperref[fam:src:D.1.4]{\ref*{fam:src:D.1.4}},
\hyperref[fam:src:D.1.6]{\ref*{fam:src:D.1.6}}).
\end{enumerate}
\hyperlink{main-I-C}{Return to Theorem C}.

\hypertarget{outline-I-D}{}
\paragraph{Theorem D.}
\begin{enumerate}[label=\arabic*.]
\item\hypertarget{step-I-D-1}{} Root characters and the cusp
normalization compute the logarithmic direct images.  The residue
connecting map is evaluated as an actual nonzero scalar, fixing
the remaining Hodge integer at one
(\hyperref[inv:A2.2]{\ref*{inv:A2.2}},
\hyperref[inv:A2.4]{\ref*{inv:A2.4}}).
\item\hypertarget{step-I-D-2}{} Coherent Leray on the curve,
the structure-sheaf inputs and Serre duality give the complete
Dolbeault table, uniformly for every admissible parameter
(\hyperref[inv:A1.0]{\ref*{inv:A1.0}},
\hyperref[inv:A3.1]{\ref*{inv:A3.1}}).
\item\hypertarget{step-I-D-3}{} The middle rows and the
$S^6$ de Rham groups force the four rank-one differentials.
The double-complex decomposition then gives the Bott--Chern and
Aeppli tables and the explicitly checked square-zero products
(\hyperref[inv:C1.1]{\ref*{inv:C1.1}},
\hyperref[inv:C2.1]{\ref*{inv:C2.1}},
\hyperref[inv:C3.3]{\ref*{inv:C3.3}}).
\item\hypertarget{step-I-D-4}{} Quotienting by $dt$ gives the
degree-one modules.  The maps $(a,b,c)\mapsto(xa-zc,yb-zc)$
and $\Omega^2\xrightarrow{\wedge dt}I$ compute their kernels
and torsion exactly; finite etale descent gives the global
support assertion
(\hyperref[inv:A4.1]{\ref*{inv:A4.1}},
\hyperref[inv:A4.2]{\ref*{inv:A4.2}}).
\end{enumerate}
\hyperlink{main-I-D}{Return to Theorem D}.

\hypertarget{outline-I-E}{}
\paragraph{Theorem E.}
\begin{enumerate}[label=\arabic*.]
\item\hypertarget{step-I-E-1}{} The primitive Picard
normalization identifies the integer twists, and weighted
coarse degrees give their section counts and four dimensions
(\hyperref[inv:C3.2]{\ref*{inv:C3.2}},
\hyperref[inv:D4.1]{\ref*{inv:D4.1}}).
\item\hypertarget{step-I-E-2}{} The actual integral cocycle
modulo the two period rows gives the regular resonance equation
and both finite root lattices.  At the cusp the normalization
complex gives exactly the reduced degree-two skyscraper
(\hyperref[inv:D4.2]{\ref*{inv:D4.2}},
\hyperref[inv:D4.3]{\ref*{inv:D4.3}}).
\item\hypertarget{step-I-E-3}{} The joint order-three
coefficient $123ik/(64\pi)$ fixes the local normal form;
Smith modules and every stabilizer character give the complete
finite stalk lengths.  Leray then gives the formula for $J_z$
(\hyperref[inv:D4.4]{\ref*{inv:D4.4}},
\hyperref[inv:D4.5]{\ref*{inv:D4.5}}).
\item\hypertarget{step-I-E-4}{} The established
$\bar\partial$-closed form $G_\gamma=G_0/12$ defines the
Picard family by $\bar\partial-2\pi izG_\gamma$.
Proper coherent base change and a real nonintegral nonresonant
point make every actual jump locus locally discrete
(\hyperref[inv:C3.1]{\ref*{inv:C3.1}},
\hyperref[inv:C3.2]{\ref*{inv:C3.2}},
\hyperref[inv:D4.5]{\ref*{inv:D4.5}}).
\item\hypertarget{step-I-E-5}{} The root-forgetting map
$(x,y)\mapsto(x^l,y^k)$ with character
$\lambda\mapsto\lambda^{kl}$ pulls back the standard Hodge
line.  Its weighted direct images and logarithmic tangent
sequence give $(1,1,1,0)$ for every stated higher sphere member
(\hyperref[push:src:E.1.2]{\ref*{push:src:E.1.2}},
\hyperref[push:src:E.2.1]{\ref*{push:src:E.2.1}}).
\item\hypertarget{step-I-E-6}{} Differentiating the actual
period groupoid gives the nonzero class $e\otimes[\gamma]$.
The induced parameter map forces the one-variable Kuranishi
ring to be reduced and smooth.  The integral translation sheaf
gives the standard automorphism group
(\hyperref[fam:st:A.5.2]{\ref*{fam:st:A.5.2}},
\hyperref[fam:src:B.4.1]{\ref*{fam:src:B.4.1}},
\hyperref[push:src:E.2.1]{\ref*{push:src:E.2.1}}).
\end{enumerate}
\hyperlink{main-I-E}{Return to Theorem E}.

\hypertarget{outline-I-F}{}
\paragraph{Theorem F.}
\begin{enumerate}[label=\arabic*.]
\item\hypertarget{step-I-F-1}{} The node charts
$X_i=W/r^i$, $Y_i=r^{i+1}/W$, $X_iY_i=r$ and their inverse
overlaps give the proper rank-one cycle.  The finite affine
actions rotate its components freely and produce the stated
cover.  Subtraction of the actual two origins fixes both deck
differences
(\hyperref[night:P2.2]{\ref*{night:P2.2}},
\hyperref[night:P2.3]{\ref*{night:P2.3}}).
\item\hypertarget{step-I-F-2}{} Resolve a proposed filling's
dominating component.  Its relative anticanonical direct image
recovers the pencil, and the three labelled critical-value
sections recover the marked base by their meromorphic cross-ratio
(\hyperref[night:P2.4]{\ref*{night:P2.4}}).
\item\hypertarget{step-I-F-3}{} Modulo $\widehat\delta$ the
actual cross-cocycles vanish.  The three unipotent annihilator
equations have determinant 6; Stein factorization on a horizontal
central component then kills its normal-meridian class
(\hyperref[night:P2.5]{\ref*{night:P2.5}},
\hyperref[night:P2.6]{\ref*{night:P2.6}}).
\item\hypertarget{step-I-F-4}{} Principalization makes both
axes normal crossing.  The Stein finite factor is birational
over the normal bidisc, so the compact corner fibre is connected.
The lifts $Q_o=T_{-mb\widehat\gamma}Q$, $R_o=R$ commute;
local meridians therefore satisfy the constant-ratio equation
(\hyperref[night:P2.7]{\ref*{night:P2.7}}).
\item\hypertarget{step-I-F-5}{} The connected incidence graph
joins a horizontal ratio zero to the strict cusp ratio $a$.
Admissibility gives $\gcd(a,12)=1$, producing the contradiction
for every proper analytic filling and every finite cover
(\hyperref[night:P2.7]{\ref*{night:P2.7}}).
\item\hypertarget{step-I-F-6}{} For higher multiplicities,
$K_X^{-N}=f^*\mathcal O(4l+3k-kl)$ with nonzero degree.
The positive canonical power and its three labelled critical
points recover the degree-one base parameter.  The same ratio
argument and $\gcd(a,N)=1$ give the full higher conclusion
(\hyperref[night:P2.9]{\ref*{night:P2.9}}).
\item\hypertarget{step-I-F-7}{} With zero logarithmic
translations the two line-bundle neck equations have zero
corner winding and give the proper flat plumbing after its
specified modification and finite quotient.  Normalization and
the double-locus smoothing formula give virtual Euler zero for
Engel's particular one-component proposal, excluding its
smoothing to $e=2$
(\hyperref[night:P2.10]{\ref*{night:P2.10}},
\hyperref[night:P2.11]{\ref*{night:P2.11}}).
\end{enumerate}
\hyperlink{main-I-F}{Return to Theorem F}.

\hypertarget{outline-I-G}{}
\paragraph{Theorem G.}
\begin{enumerate}[label=\arabic*.]
\item\hypertarget{step-I-G-1}{} The canonical divisor recovers
$S_2$ intrinsically in a proper extension.  Its normalization
and Albanese invariant retain $J(iw/216)$, whose natural boundary
excludes crossing the outer circle
(\hyperref[push:src:D.1.3]{\ref*{push:src:D.1.3}}).
\item\hypertarget{step-I-G-2}{} The fixed complex-linear and
integral basis changes give the displayed four periods.
The local theta expansion gives $B(h)=-3ih^2+O(h^4)$ and
the exact inverse transpose gives every integral dual covector
(\hyperref[push:src:D.2.1]{\ref*{push:src:D.2.1}},
\hyperref[eq:d-dual-exact]{\ref*{eq:d-dual-exact}}).
\item\hypertarget{step-I-G-3}{} Bounded dual covectors have
constant integral labels after subselection; their labels form
$L_a$.  The quadratic real term cancels, leaving
$\lambda\cdot V_n+o(1)$.  The primitive surviving subgroup and
its limiting homomorphism give exactly $\Gamma_\infty$
(\hyperref[push:src:D.3.1]{\ref*{push:src:D.3.1}}).
\item\hypertarget{step-I-G-4}{} Choosing successively separated
scales realizes every primitive subgroup and limiting
homomorphism.  Pontryagin--Chabauty duality converts their dual
graphs to the stated real-torus additive quotients
(\hyperref[push:src:D.3.1]{\ref*{push:src:D.3.1}}).
\item\hypertarget{step-I-G-5}{} At fixed boundary parameter,
irrational root rays make all noncentral dual labels escape.
The surviving rational label is
$\operatorname{lcm}(2,n)\Re\xi$; its circle-valued function
separates precisely the asserted classes.  Collapsing those
classes proves the universal Hausdorff-reflection property
(\hyperref[push:src:D.2.3]{\ref*{push:src:D.2.3}}).
\end{enumerate}
\hyperlink{main-I-G}{Return to Theorem G}.

% RungZero blocks 1 and 25 are deliberately assigned to the shared
% series overview and the final open-problems part, respectively.
% The literal coefficient-atlas question is inv:C3.8; the arithmetic
% question is the retained ari:C.7.1--ari:C.7.12 discussion, with the
% explicit unresolved relative-period gate at ari:C.7.4.

\end{document}
